\begin{statementbox}
\textbf{焦准比例。} 双曲线 $x^2/a^2-y^2/b^2=1$，$a,b>0$，
$c=\sqrt{a^2+b^2}>a$，$e=c/a>1$。
$F_R(c,0)$ 对应 $\ell_R:x=a^2/c$；$F_L(-c,0)$ 对应 $\ell_L:x=-a^2/c$。
对双曲线两支上的每一点 $P$，
\[
\boxed{\frac{PF_R}{d(P,\ell_R)}=\frac{PF_L}{d(P,\ell_L)}=e}.
\]
准线位于两顶点之间，因为 $0<a^2/c<a$；而双曲线上 $|x|\ge a$，准距均不为零。
右焦半径可统一写为 $\boxed{PF_R=|ex-a|}$，右准距为 $|x-a^2/c|$。

\textbf{轨迹的反向。} 给定 $F\notin\ell$ 和 $e>1$，点集 $PF=e\,d(P,\ell)$ 是整条非退化双曲线。
同一焦准对适用于两支。焦点在 $y$ 轴时交换坐标，准线为 $y=\pm a^2/c$。
\end{statementbox}
\clearpage
\begin{statementbox}

\medskip
\textbf{一句直观话：} \emph{离焦点有多远，恰好是离对应准线有多远的 $e$ 倍。}
其中 $d(P,\ell)$ 指\textbf{最短的垂直距离}，不是从 $P$ 到准线上随意选一点的距离。

\textbf{把“两套焦准”写完整。}
\[
\begin{array}{ll}
\text{右焦准对：}&F_R(c,0),\quad \ell_R:x=\dfrac{a^2}{c},\\[3pt]
\text{左焦准对：}&F_L(-c,0),\quad \ell_L:x=-\dfrac{a^2}{c}.
\end{array}
\]
对应的两条\textbf{全双曲线}焦半径公式为
\[
\boxed{PF_R=\left|\frac ca x-a\right|},\qquad
\boxed{PF_L=\left|\frac ca x+a\right|}.
\]
这两个等式不要求点 $P$ 与焦点在同一支：\textbf{焦点与准线要配对，点 $P$ 可以在任意一支。}

\textbf{纵向情形。} 若 $\dfrac{y^2}{a^2}-\dfrac{x^2}{b^2}=1$，则
$F_U(0,c)$ 配 $\ell_U:y=a^2/c$，$F_D(0,-c)$ 配 $\ell_D:y=-a^2/c$，且
\[
\boxed{PF_U=\left|\frac ca y-a\right|
=e\,d(P,\ell_U)},\qquad
\boxed{PF_D=\left|\frac ca y+a\right|
=e\,d(P,\ell_D)}.
\]
\textbf{注意前提：} 以上公式中的 $a$ 是实半轴长，$c$ 是焦距之半，
横向或纵向均有 $c^2=a^2+b^2$；一般焦准轨迹还要求焦点 $F\notin\ell$。

\end{statementbox}
